\section{Quantum algorithms}

We now review standard Quantum algorithms in the tensor network formalism.

\subsection{Quantum Fourier transform}\glsadd{cc:FT:QFT}

Recall the tensor network decomposition of the Fourier transform developed in \secref{sec:fourierTransform}.

When $\vectorat{\shortcatvariables}$ is a quantum state, the Fourier transform can be understood as a unitary evolving the state.
Since both the controlled $U$ tensor and the Hadamard $H$ are unitary, each of them can be implemented by a sequence of gates.
For sparse implementation one can exploit that $U$ is a tensor product of matrices (i.e. single qubit gates).

% Gate representation
We introduce the gates
\begin{align*}
    R^{\catenumerator}[\ccatvariable{\insymbol},\ccatvariable{\outsymbol}] =
    \coloredmatrixof{
        1 & 0 \\
        0 & \expof{2\pi i \frac{1}{2^\catenumerator}}
    }{\ccatvariable{\insymbol},\ccatvariable{\outsymbol}} \, .
\end{align*}
The unitary $U^{\catorder}$ in the tensor network decomposition of the fourier transform (see \secref{sec:ft:tn}) is then the tensor product
\begin{align*}
    U^{\catorder}[\ccatvariableof{\insymbol}{[\catorder]},\ccatvariableof{\outsymbol}{[\catorder]}]
    = \bigotimes_{\catenumeratorin} R^{\catenumerator+2}[\ccatvariableof{\insymbol}{\catenumerator},\ccatvariableof{\outsymbol}{\catenumerator}] \, .
\end{align*}
Controlling $U^{\catorder}$ by $\catvariableof{\catorder}$ is equivalent to controlling each $R^{\catenumerator}$ gate.


% Gate complexity
To realize the QFT we need $\catorder$ single-qubit Hadamard gates and $\sum_{\catenumeratorin} \catenumerator = \frac{\catorder\cdot(\catorder-1)}{2}$ controlled diagonal single-qubit gates (to realize the $U^{\catenumerator}$ transforms).
The gate complexity is therefore $\mathcal{O}(\catorder^2)$.

\subsubsection{Walsh-Hadamard transform}\glsadd{qc:WH}

In case of $\catdimof{\catenumerator}=2$, the Fourier transform is known to be the Walsh-Hadamard tranform.
The transform is performed by applying Hadamard gates on each variable.
This is equivalent with contraction of the tensor
\begin{align*}
    \hgate^{\catorder}\left[\catvariableof{[\catorder],\insymbol},\catvariableof{[\catorder],\outsymbol}\right] =
    \bigotimes_{\catenumeratorin} \hgateat{\ccatvariableof{\catenumerator}{\insymbol},\ccatvariableof{\catenumerator}{\outsymbol}} \, ,
\end{align*}
which coordinates are
\begin{align*}
    \hgate^{\catorder}\left[\indexedcatvariableof{[\catorder],\insymbol},\indexedcatvariableof{[\catorder],\outsymbol}\right] =
    \sqrt{\frac{1}{2^{\catorder}}} \left(-1\right)^{\sum_{\catenumeratorin}\catindexof{\catenumerator,\insymbol} \cdot \catindexof{\catenumerator,\outsymbol}} \, .
\end{align*}

We have
\begin{align*}
    \contractionof{\onehotmapofat{0}{\ccatvariable{\insymbol}},
        \hgateat{\ccatvariable{\insymbol},\ccatvariable{\outsymbol}}}{\ccatvariable{\outsymbol}}
    &= \sqrt{\frac{1}{2}} \onesat{\ccatvariable{\outsymbol}} \\
    \contractionof{\onehotmapofat{1}{\ccatvariable{\insymbol}},
        \hgateat{\ccatvariable{\insymbol},\ccatvariable{\outsymbol}}}{\ccatvariable{\outsymbol}}
    &= \sqrt{\frac{1}{2}} \coloredmatrixof{1 \\ -1}{\ccatvariable{\outsymbol}}
\end{align*}
and conversely
\begin{align*}
    \contractionof{\sqrt{\frac{1}{2}} \onesat{\ccatvariable{\insymbol}},
        \hgateat{\ccatvariable{\insymbol},\ccatvariable{\outsymbol}}}{\ccatvariable{\outsymbol}}
    &= \onehotmapofat{0}{\ccatvariable{\outsymbol}} \\
    \contractionof{\sqrt{\frac{1}{2}} \coloredmatrixof{1 \\ -1}{\ccatvariable{\insymbol}},
        \hgateat{\ccatvariable{\insymbol},\ccatvariable{\outsymbol}}}{\ccatvariable{\outsymbol}}
    &= \onehotmapofat{1}{\ccatvariable{\outsymbol}} \, .
\end{align*}
If and only if the function is constant, then the transformed state is the ground state.

Let us present two known algorithms as the measurement of states prepared by \computationCircuits{}.

\begin{example}[Deutsch-Jozsa Test]
    We measure the probability of the ground state of the sign encoding after a Walsh-Hadamard transform.
    The probability of measuring the ground state after a Walsh-Hadamard transform of the distributed qubits is then
    \begin{align*}
        \frac{1}{2^{\catorder}} \cdot \contraction{\absof{\contractionof{\signqstateofat{\exformula}{\shortcatvariables,\headvariable}}{\headvariable}}^2}
        = \frac{1}{2^{2\cdot\catorder}} \left(\contraction{\exformula}-\contraction{\lnot\exformula}\right)^2 \, .
    \end{align*}
    Thus, if and only if the probability of the ground state is $1$, we have that the formula is constant, that is $\exformula\in\{\top,\bot\}$.
    If and only if the probability of the ground state is $0$, we have that the formula is balanced, that is $\contraction{\exformula}=\contraction{\lnot\exformula}$.
    When its known, that the formula is either constant or balanced, we can decide which one it is by measuring the probability of the ground state.
    This is the Deutsch-Josza algorithm \cite{deutsch_rapid_1992}.

    %%
    Note that when allowing for an exponentially small error the decision can be efficiently performed by a classical algorithm.
    When drawing $n$ samples and deciding on $\meanparam=0.5$ if and only if in the samples a $0$ and a $1$ is observed, we have an error probability of $0$ when $\meanparam\in\{0,1\}$ and of
    \begin{align*}
        \probof{E|\meanparam=0.5} = 2 \cdot \frac{1}{2^{n}} \, .
    \end{align*}
    The quantum advantage is thus bound to the demand of an error-free decision, for which a generic classical algorithm would iterate through all states and is therefore exponentially slower than the quantum alternative.

    \begin{figure}
        \begin{center}

            \begin{tikzpicture}[scale=0.35,thick]

                \draw (0,0) rectangle (2,7.5);
                \node[anchor=center] (text) at (1,3.75) {\corelabelsize $\signqstateof{\exformula}$};

                \draw (2,6.5) -- (4,6.5) node[midway,above] {\colorlabelsize $\headvariable$};
                \drawqcmeasuresymbol{5}{6.5}
                \draw (6,6.5) -- (8,6.5); % node[midway,above] {\colorlabelsize $\headvariable$};
                \drawvariabledot{8}{6.5}

                \draw (2,4) -- (4,4) node[midway,above] {\colorlabelsize $\catvariableof{0}$};
                \draw (4,3) rectangle (6,5);
                \node[anchor=center] at (5,4) {$\hgate$};
                \draw (6,4) -- (7,4);
                \drawqcmeasuresymbol{8}{4}
                \draw (9,4) -- (10,4);
                \draw (10,3) rectangle (12,5);
                \node[anchor=center] at (11,4) {$\fbasis$};

                \node[anchor=center] (text) at (3,3.25) {$\vdots$};
                \draw (2,1) -- (4,1) node[midway,above] {\colorlabelsize $\catvariableof{\catorder\shortminus 1}$};
                \draw (4,0) rectangle (6,2);
                \node[anchor=center] at (5,1) {$\hgate$};
                \draw (6,1) -- (7,1);
                \drawqcmeasuresymbol{8}{1}
                \draw (9,1) -- (10,1);
                \draw (10,0) rectangle (12,2);
                \node[anchor=center] at (11,1) {$\fbasis$};

                \node[anchor=center] (text) at (14.5,3.75) {${=}\,\, \frac{1}{2^{2\cdot\catorder}}$};

                \begin{scope}[shift={(18,1.5)}]

                    % Squared
                    \draw (-1,6) to[bend right=15] (-1,-2);
                    \draw (7,6) to[bend right=-15] (7,-2);
                    \drawtextnode{7.5}{6.5}{$2$}

                    \draw (1.5,3) rectangle (4.5,6);
                    \drawtextnode{3}{4.5}{
                        $\begin{bmatrix}
                             -1 \\
                             1
                        \end{bmatrix}$
                    }

                    \draw[->-] (3,1) -- (3,3) node[midway,right] {\colorlabelsize $\headvariable$};
                    \draw (0,-1) rectangle (6,1);
                    \drawtextnode{3}{0}{\corelabelsize $\bencodingof{\exformula}$}
                    \draw[->-] (1,-2) -- (1,-1) node[midway,left] {\colorlabelsize $\catvariableof{0}$};
                    \drawvariabledot{1}{-2}
                    \drawtextnode{3}{-1.5}{$\cdots$}
                    \draw[->-] (5,-2) -- (5,-1) node[midway,right] {\colorlabelsize $\catvariableof{\catorder\shortminus1}$};
                    \drawvariabledot{5}{-2}
                \end{scope}

            \end{tikzpicture}
        \end{center}
        \caption{Measurement setup for the Deutsch-Jozsa algorithm on the formula $\exformula$.
        }\label{fig:deutschJozsa}
    \end{figure}

\end{example}


% Generalization
\todo{Another generalizations would be period detection.}
Deutsch-Jozsa can be generalized to estimating the moments of a distribution given its q-sample.
Note, that we have to demand a q-sample (i.e. vanishing phase tensor), otherwise the measurement is affected by the phase tensor and the Deutsch-Jozsa test is effectively done wrt another function.
If for example the coordinates of the phase tensor are multiples of $\pi$, and the phases are thus $+1/-1$, then the test is done wrt the function $\exformula\oplus(\frac{1}{\pi}\phasecore)$ instead of $\exformula$.

\begin{example}[Inversion Test]
    Concatenating the \computationCircuits{} to two formulas $\exformula$ and $\secexformula$ prepares the state
    \begin{align*}
        \basqstateofat{\exformula\oplus\secexformula}{\shortcatvariables,\headvariable}
    \end{align*}
    The probability of measuring the ground state after a Walsh-Hadamard transform of the distributed qubits (see \figref{fig:inversionTest}) is then
    \begin{align*}
        \frac{1}{2^{\catorder}} \absof{\contractionof{\basqstateofat{\exformula\oplus\secexformula}{\shortcatvariables,\headvariable}}{\headvariable}}^2
        &= \frac{1}{2^{2\cdot\catorder}}
        \begin{bmatrix}
            \contraction{\exformula\oplus\secexformula}
            \contraction{\lnot(\exformula\oplus\secexformula)}
        \end{bmatrix} \\
        &= \frac{1}{2^{2\cdot\catorder}}
        \begin{bmatrix}
            \#{\left\{\shortcatindices\in\atomstates \wcols \exformulaat{\indexedshortcatvariables}=\secexformula\left[\indexedshortcatvariables\right]\right\}} \\
            \#{\left\{\shortcatindices\in\atomstates \wcols \exformulaat{\indexedshortcatvariables}\neq\secexformula\left[\indexedshortcatvariables\right]\right\}}
        \end{bmatrix} \, .
    \end{align*}
    The second equation holds since $\lnot(\exformula\oplus\secexformula)$ is equivalent with the biconditional $\exformula\Leftrightarrow\secexformula$.

    % Inversion test
    This is in more generality the inversion test of measuring the overlap of $\basqstateof{\exformula}$ with $\basqstateof{\secexformula}$ (see \cite{schuld_machine_2021}).
    Here we used that the \computationCircuit{} is self-adjoint and understand the concatenation as the unitaries preparing $\basqstateof{\exformula}$ and the adjoint (together with the Walsh-Hadamard transform) of $\basqstateof{\secexformula}$.

    \begin{figure}
        \begin{center}

            \begin{tikzpicture}[scale=0.35,thick]

                \draw (0,0) rectangle (2,7.5);
                \node[anchor=center] (text) at (1,3.75) {\corelabelsize $\basqstateof{\exformula\oplus\secexformula}$};

                \draw (2,6.5) -- (4,6.5) node[midway,above] {\colorlabelsize $\headvariable$};
                \drawqcmeasuresymbol{5}{6.5}
                \draw (6,6.5) -- (8,6.5) node[midway,above] {\colorlabelsize $\headvariable$};

                \draw (2,4) -- (4,4) node[midway,above] {\colorlabelsize $\catvariableof{0}$};
                \draw (4,3) rectangle (6,5);
                \node[anchor=center] at (5,4) {$\hgate$};
                \draw (6,4) -- (7,4);
                \drawqcmeasuresymbol{8}{4}
                \draw (9,4) -- (10,4);
                \draw (10,3) rectangle (12,5);
                \node[anchor=center] at (11,4) {$\fbasis$};

                \node[anchor=center] (text) at (3,3.25) {$\vdots$};
                \draw (2,1) -- (4,1) node[midway,above] {\colorlabelsize $\catvariableof{\catorder\shortminus 1}$};
                \draw (4,0) rectangle (6,2);
                \node[anchor=center] at (5,1) {$\hgate$};
                \draw (6,1) -- (7,1);
                \drawqcmeasuresymbol{8}{1}
                \draw (9,1) -- (10,1);
                \draw (10,0) rectangle (12,2);
                \node[anchor=center] at (11,1) {$\fbasis$};

                \node[anchor=center] (text) at (14.5,3.75) {${=}\,\, \frac{1}{2^{2\cdot\catorder}}$};

                \begin{scope}[shift={(18,3.75)}]

                    \draw (3,-0.5) ellipse (4.5cm and 3cm);
                    \node[anchor=center] at (5,3.25) {$\absof{\cdot}^2$};

                    \draw[->-] (3,1) -- (3,3) node[midway,right] {\colorlabelsize $\headvariable$};
                    \draw (0,-1) rectangle (6,1);
                    \drawtextnode{3}{0}{\corelabelsize $\bencodingof{\exformula\oplus\secexformula}$}
                    \draw[->-] (1,-2) -- (1,-1) node[midway,left] {\colorlabelsize $\catvariableof{0}$};
                    \drawvariabledot{1}{-2}
                    \drawtextnode{3}{-1.5}{$\cdots$}
                    \draw[->-] (5,-2) -- (5,-1) node[midway,right] {\colorlabelsize $\catvariableof{\catorder\shortminus1}$};
                    \drawvariabledot{5}{-2}
                \end{scope}

            \end{tikzpicture}
        \end{center}
        \caption{Measurement setup for the inversion test on the formulas $\exformula$ and $\secexformula$.
        We measure the probability of the statistic qubit $\headvariable$, when the distributed qubits $\shortcatvariables$ are in the ground state after a Walsh-Hadamard transform.
        This is equal to a scaled and squared contraction of the basis encoding to the formula $\exformula\oplus\secexformula$.
        }\label{fig:inversionTest}
    \end{figure}
\end{example}


%\subsection{Further Overlap-measuring Circuits}
Quantum Circuits such as the SWAP test and the Hadamard test (\cite[Section 3.6.1]{schuld_machine_2021}) can be used to measure overlaps of quantum states, which are the squared absolutes of contractions of two state tensors.
\begin{itemize}
    \item When we have preparation schemes for two tensors, we can control them with a common ancilla qubit and measure their contraction by a Hadamard test (alternatively, using the SWAP test and state preparation in two registers).
    \item Can we extend these schemes to contractions of more general tensor networks?
\end{itemize}

%\subsection{Discussion}
While the preparation of contracted basis encodings is done efficiently, accessing the prepared state by measurements is a bottleneck.
This is a typical \emph{quantum state tomography} problem \cite{roth_semi-device-dependent_2023}.
Sparsity of the contracted state can be exploited to probabilistically guarantee precise estimates \cite{gross_quantum_2010}.
However, when deciding whether $\exformula$ is satisfiable, based on the prepared contraction $\sqrt{\normalizationof{\bencodingofat{\exformula}{\headvariable,\shortcatvariables}}{\headvariable}}$, we have an exponentially small error tolerance.
Moreover, the applied sample mean is the UMVUE (Uniformly Minimum Variance Unbiased Estimator, see \cite{casella_statistical_2001}), and can thus not be improved by other unbiased estimators.




\subsubsection{Period detection}

QFT is a typical component in Quantum algorithms: QFT can be done efficiently with $\mathcal{0}(\catorder^2)$ controlled single-qubit gates.
If the Fourier transform is a basis vector, a single measurement is sufficient to determine it.
We now investigate how this can be used in the detection of periods.

%\subsubsection{Detecting periods}

Based on the support of the Fourier transform, one can derive the periods of a function.
Periods are indices $(\thirdcatvariableof{0},\ldots,\thirdcatvariableof{\catorder-1})$ such that
\begin{align*}
    \hypercoreat{\shortcatvariables=\shortcatindices}
    = \hypercoreat{\shortcatvariables=\shortcatindices\oplus\thirdcatindexof{[\catorder]}}
\end{align*}
where $\oplus$ is the position-wise sum $\mathrm{mod}\,\,\catdimof{\catenumerator}$.


Each supported index of the Fourier transform corresponds with periods in the function.
The periods of the functions are thus at least the intersection of those (when assuming more structure one might have more periods).

%% WH-case
In case of the Walsh-Hadamard transform (i.e. $\catdimof{\catenumerator}=2$), we can easily determine the periods corresponding to each basis vector.
The inverse Fourier transform of $\onehotmapofat{\selindexof{[\catorder]}}{\selvariableof{[\catorder]}}$ is the tensor
\begin{align*}
    \left(\bigotimes_{\catenumeratorin \wcols \selindexof{\catenumerator}=0} \hplusbasat{\catvariableof{\catenumerator}} \right)
    \otimes
    \left(\bigotimes_{\catenumeratorin \wcols \selindexof{\catenumerator}=1} \hminusbasat{\catvariableof{\catenumerator}} \right) \, .
\end{align*}

For each $\hplusbasat{\catvariableof{\catenumerator}}$ we have a period generator of $\onehotmapof{\catenumerator}$.
For each pair of $\hminusbasat{\catvariableof{\catenumerator}}\otimes \hminusbasat{\catvariableof{\seccatenumerator}}$ we have a period generator of $\onehotmapof{\catenumerator}+\onehotmapof{\seccatenumerator}$.

One example is the Deutsch-Jozsa test, where the transformed is the $\onehotmapofat{0}{\selvariableof{[\catorder]}}$ state, if and only if the function is constant.
In this case, any vector is a period, since the group itself is generated by the period generators.


%\subsubsection{Generic Hidden Subgroup Problems}

% Generic hidden subgroup problems
%In general \cite{nielsen_quantum_2010} approached by generic Fourier transforms of the distributed variables of a basis encoding state to a function.

\subsection{Quantum Phase Estimation}\label{sec:quantumPhaseEstimation}\glsadd{qc:QPE}

% Generic Quantum Phase Estimation
When having a circuit implementing a unitary, the phase of the eigenvalues can be estimated using the Quantum Phase Estimation routine \cite{kitaev_quantum_1995, cleve_quantum_1998}.

Let $\unitaryat{\ccatvariableof{\insymbol}{[\catorder]},\ccatvariableof{\outsymbol}{[\catorder]}}$ be a unitary then the absolute of any eigenvalue is $1$ and the eigenvalue is expressible by
\begin{align*}
    \expof{\twopii \cdot \phi}
\end{align*}
for some $\phi\in[0,1]$.
Let $\phi=0.\phi_{0}\cdots\phi_{n-1}$ be a $2$-adic representation of $\phi$, that is
\begin{align*}
    \phi
    = \sum_{\catenumerator\in[n]} 2^{-\catenumerator-1} \phi_{\catenumerator} \, .
\end{align*}

In the Quantum Phase Estimation routine the Fourier transform of $\onehotmapofat{\phi_0,...,\phi_{n-1}}{\selvariableof{[n]}}$
\begin{align*}
    \sqrt{\frac{1}{2^{n}}} \cdot \bigotimes_{\catenumerator\in[n]} \left(\fbasisat{\selvariableof{0}} + \expof{\twopii \cdot 0.\phi_{n-\catenumerator-1}\cdots\phi_{n-1}}\right)
\end{align*}
is prepared as explained below.
The inverse Fourier transform on the register $\selvariableof{[n]}$ results then in the state
\begin{align*}
    \onehotmapofat{\phi_0,...,\phi_{n-1}}{\selvariableof{[n]}} \otimes \qstateat{\shortcatvariables}
\end{align*}
and $\phi_{\catenumerator}$ for $\catenumerator\in[n]$ can be retrieved by a single computational basis measurement of the register $\selvariableof{[n]}$.

To prepare the Fourier transform of $\onehotmapofat{\phi_0,...,\phi_{n-1}}{\shortcatvariables}$, we start with the state
\begin{align*}
    \hplusbasat{\selvariable} \otimes \qstateat{\shortcatvariables} \, .
\end{align*}
When applying $\unitaryat{\ccatvariableof{\insymbol}{[\catorder]},\ccatvariableof{\outsymbol}{[\catorder]}}$ $2^{\catenumerator}$ times controlled on $\selvariable$ we get the state
\begin{align*}
    \sqrt{\frac{1}{2}} \left(\fbasisat{\selvariableof{0}} + \expof{\twopii \cdot 0.\phi_{n-\catenumerator-1}\cdots\phi_{n-1}}\right) \otimes \qstateat{\shortcatvariables}
\end{align*}

Now we apply for each $\catenumerator\in[n]$ the unitary $\unitaryat{\ccatvariableof{\insymbol}{[\catorder]},\ccatvariableof{\outsymbol}{[\catorder]}}$ $2^{\catenumerator}$ times controlled on $\selvariableof{\catenumerator}$ and get
\begin{align*}
    \sqrt{\frac{1}{2^{n}}} \onesat{\selvariableof{[n]}} \otimes \qstateat{\shortcatvariables} \, .
\end{align*}

\subsection{Grover operator}\label{sec:groverOperator}

We derive quantum inference schemes based on the Grover operator, which we want to first introduce in full generality.

\subsubsection{Definition}

Let us assume we have two orthogonal quantum states $\qstateofat{\goodstatesymbol}{\shortcatvariables},\qstateofat{\badstatesymbol}{\shortcatvariables}$ spanning a subspace $\subspaceof{}=\spanof{\qstateofat{\goodstatesymbol}{\shortcatvariables},\qstateofat{\goodstatesymbol}{\shortcatvariables}}$ and let us choose another state
\begin{align*}
    \qstateofat{0}{\shortcatvariables}
    \in \subspaceof{}\, .
\end{align*}
Then there is $\rotanglesymbol\in[0,4\pi)$ with
\begin{align*}
    \qstateofat{0}{\shortcatvariables}
    = \sinof{\frac{\rotanglesymbol}{2}}\qstateofat{\goodstatesymbol}{\shortcatvariables}
    + \cosof{\frac{\rotanglesymbol}{2}}\qstateofat{\badstatesymbol}{\shortcatvariables} \, .
\end{align*}
Consider the reflection about the states $\qstateof{\badstatesymbol}$ and $\qstateof{0}$ by
\begin{align*}
    \reflectionofat{\badstatesymbol}{\inshortcatvariables,\outshortcatvariables}
    &= 2 \cdot \qstateofat{\badstatesymbol}{\inshortcatvariables} \otimes \qstateofat{\badstatesymbol}{\outshortcatvariables} - \identityat{\inshortcatvariables,\outshortcatvariables} \\
    \reflectionofat{0}{\inshortcatvariables,\outshortcatvariables}
    &= 2 \cdot \qstateofat{0}{\inshortcatvariables} \otimes \qstateofat{0}{\outshortcatvariables} - \identityat{\inshortcatvariables,\outshortcatvariables} \, ,
\end{align*}
and their concatenation
\begin{align*}
    \groverat{\inshortcatvariables,\outshortcatvariables}
    = \contractionof{
        \reflectionofat{\badstatesymbol}{\inshortcatvariables,\ccatvariableof{[\catorder]}{\midsymbol}},
        \reflectionofat{0}{\ccatvariableof{[\catorder]}{\midsymbol},\outshortcatvariables}
    }{\inshortcatvariables,\outshortcatvariables}
\end{align*}
which is called the \emph{Grover operator}.

\subsubsection{Action on projected quantum states}

The Grover operator leaves any quantum state $\qstate$ invariant, which is orthogonal to $\qstateof{\goodstatesymbol}$ and $\qstateof{\badstatesymbol}$ (such states get a factor $(-1)$ by both reflection operations).
The projection of the quantum state onto $\subspaceof{}$ gets rotated by the Grover operator (see \figref{fig:rotationSketch}).
To be more precise, we denote the projection by $(\mathcal{P}\qstate)[\selvariable]$ which has the coordinates
\begin{align*}
    \left(\mathcal{P}_{\subspaceof{}}\qstate\right)\left[\selvariable=0\right]
    &= \contraction{\left(\qstatewith\right)^{\dagger},\qstateofat{\goodstatesymbol}{\shortcatvariables}}\\
    \left(\mathcal{P}_{\subspaceof{}}\qstate\right)\left[\selvariable=1\right]
    &= \contraction{\left(\qstatewith\right)^{\dagger},\qstateofat{\goodstatesymbol}{\shortcatvariables}} \, .
\end{align*}
We define a matrix representation of the Grover operator as
\begin{align*}
    \groverofat{\mathrm{red}}{\cselvariable{\outsymbol},\cselvariable{\insymbol}}
    = \coloredmatrixof{
        \cosof{\rotanglesymbol} & -\sinof{\rotanglesymbol} \\
        \sinof{\rotanglesymbol} & \cosof{\rotanglesymbol}
    }{\cselvariable{\outsymbol},\cselvariable{\insymbol}}
\end{align*}
and it holds for any quantum state $\qstatewith$ that
\begin{align*}
    \left(\mathcal{P}_{\subspaceof{}}\groversymbol\qstate\right)\left[\cselvariable{\outsymbol}\right]
    = \contractionof{\groverofat{\mathrm{red}}{\cselvariable{\outsymbol},\cselvariable{\insymbol}},
        \left(\mathcal{P}_{\subspaceof{}}\qstate\right)\left[\cselvariable{\insymbol}\right]
    }{\outsymbol} \, .
\end{align*}
The Grover operator therefore has the eigenvalues:
\begin{itemize}
    \item $\expof{\pm i \rotanglesymbol}$ with the eigen spaces spanned by $\qstateof{\goodstatesymbol}\pm i \qstateof{\badstatesymbol}$
    \item $1$ with the eigen space by the orthogonal complement of $\subspaceof{}$
\end{itemize}
Application the Grover operator $\repnum\in\nn$ times on $\qstateof{0}$ results in the state
\begin{align*}
    \qstateofat{\repnum}{\shortcatvariables}
    = \sinof{\left(\frac{1}{2}+\repnum\right)\rotanglesymbol} \qstateofat{\goodstatesymbol}{\shortcatvariables}
    + \cosof{\left(\frac{1}{2}+\repnum\right)\rotanglesymbol} \qstateofat{\badstatesymbol}{\shortcatvariables} \, .
\end{align*}

\begin{figure}[t]
    \begin{center}
        \begin{tikzpicture}[>=stealth, scale=4]

    % Coordinates
    \coordinate (O) at (0,0);
    \coordinate (W) at (0,1);     % Target state |w>
    \coordinate (Sprime) at (1,0); % All other states |s'>

    % Draw Axes
    \draw[->, thick] (-0.1,0) -- (1.1,0) node[right] {\corelabelsize $\qstateof{\badstatesymbol}$};
    \draw[->, thick] (0,-0.1) -- (0,1.1) node[above] {\corelabelsize $\qstateof{\goodstatesymbol}$};

    \def\sthalfangle{15}
    \def\nextangle{3*\sthalfangle} % 3 * theta
    \draw[->, thick] (O) -- (\sthalfangle:0.9) node[right] {\corelabelsize $\qstateof{0}=\sinof{\frac{\rotanglesymbol}{2}}\qstateof{\goodstatesymbol}+\cosof{\frac{\rotanglesymbol}{2}}\qstateof{\badstatesymbol}$};
    \draw[dashed, thick] (0.4,0) arc (0:\sthalfangle:0.4) node[midway, right, xshift=2pt] {\corelabelsize $\frac{\rotanglesymbol}{2}$};

    \draw[<-, thick, blue] (-1*\sthalfangle:0.6) arc (-\sthalfangle:\sthalfangle:0.6) node[midway, right, yshift=-5pt, xshift=2pt] {\corelabelsize $\rotanglesymbol$};

    % Draw the Oracle Reflection (reflected across s' axis)
    \draw[->, thick, blue] (O) -- (-\sthalfangle:0.9) node[right] {\corelabelsize $-\sinof{\frac{\rotanglesymbol}{2}}\qstateof{\goodstatesymbol}+\cosof{\frac{\rotanglesymbol}{2}}\qstateof{\badstatesymbol}$};


    \draw[->, thick, red] (-1*\sthalfangle:0.8) arc (-\sthalfangle:3*\sthalfangle:0.8) node[midway, right, yshift=15pt, xshift=-5pt] {\corelabelsize $2\rotanglesymbol$};

    % Draw the first Grover Iteration (rotated by 2*theta from original |s>)

    \draw[->, red, thick] (O) -- (\nextangle:0.9) node[above right] {\corelabelsize $\qstateof{1}=\sinof{\left(\frac{1}{2}+1\right)\rotanglesymbol}\qstateof{\goodstatesymbol}+\cosof{\left(\frac{1}{2}+1\right)\rotanglesymbol}\qstateof{\badstatesymbol}$};

    \node[below left] at (O) {\corelabelsize $0$};

\end{tikzpicture}
    \end{center}
    \caption{Grover operator as a rotation in $\spanof{\qstateof{\goodstatesymbol},\qstateof{\badstatesymbol}}$.}
    \label{fig:rotationSketch}
\end{figure}

\subsubsection{Angle measurement}\glsadd{qc:QPE:grover}\label{sec:groverAngleMeasurement}

Using \gls{qc:QPE} we estimate the angle of the Grover rotation.

The eigenvalues of the Grover operator are (see \secref{sec:groverOperator}) $\expof{\pm \mathrm{i}\rotanglesymbol}$ and $1$, and have therefore phases $\frac{\pm\rotanglesymbol}{\twopii}$ and $0$.
Note that $\spanof{\qstateof{\goodstatesymbol},\qstateof{\badstatesymbol}}$ is also spanned by the first two eigenvectors.
When the state of the register is in $\spanof{\qstateof{\goodstatesymbol},\qstateof{\badstatesymbol}}$ we therefore measure either $\frac{\rotanglesymbol}{\twopii}$ or $\frac{\twopii-\rotanglesymbol}{\twopii}$.


\subsubsection{Circuit Realization of the reflection operators}

Let us now investigate how the reflection operations in the Grover operator can be realized by circuits.

%\subsubsection{Construction based on controlled Pauli-Z}
First, the reflection about the state $\onehotmapofat{\onetuple{[\catorder]}}{\shortcatvariables}$ is realized by a multiple controlled Pauli-Z (see \figref{fig:reflectionZRealization})
%    \begin{align*}
%        -\reflectionofat{\goodstatesymbol}{\inshortcatvariables,\outshortcatvariables}
%        = \mathrm{MC}^{\seldim-1}\paulizat{\cavariableof{\seldim-1}{\insymbol},\cavariableof{\seldim-1}{\outsymbol},\avariableof{[\seldim-1]}}
%    \end{align*}
\begin{align*}
    -\reflectionofat{\onehotmapof{\onetuple{[\catorder]}}}{\inshortcatvariables,\outshortcatvariables}
    = \contractionof{
        \mathrm{MC^{\catorder-1}}\paulizat{\ccatvariableof{\catorder-1}{\insymbol},\ccatvariableof{\catorder-1}{\outsymbol},\catvariableof{[\catorder-1]}},
        \identityat{\ccatvariableof{[\catorder-1]},\ccatvariableof{[\catorder-1]}{\insymbol},\ccatvariableof{[\catorder-1]}{\outsymbol}}
    }{\inshortcatvariables,\outshortcatvariables}
\end{align*}
Note that a global phase of $-1$ is introduced, which leaves the measurement distribution invariant.
% Extending to e_1
Adding layers of Pauli-X gates to each $\catenumeratorin$ before and after $-\reflectionofat{\onehotmapof{\onetuple{[\catorder]}}}{\inshortcatvariables,\outshortcatvariables}$ (see \figref{fig:reflectionZRealization}) prepares the negated reflection $-\reflectionofat{\onehotmapof{\zerotuple{[\catorder]}}}{\inshortcatvariables,\outshortcatvariables}$, since
\begin{align*}
    \tbasisat{\ccatvariableof{\catenumerator}{\outsymbol}}
    = \contractionof{
        \paulizat{\ccatvariableof{\catenumerator}{\outsymbol},\ccatvariableof{\catenumerator}{\insymbol}},\fbasisat{\ccatvariableof{\catenumerator}{\insymbol}}
    }{\ccatvariableof{\catenumerator}{\outsymbol}} \, .
\end{align*}
% Extending to U
Having a unitary $\unitarysymbol$ preparing $\qstate$ by action on the ground state, we realize the reflection (see \figref{fig:reflectionZRealization})
\begin{align*}
    \reflectionofat{\qstate}{\inshortcatvariables,\outshortcatvariables}
    = \contractionof{
        \unitaryofat{\dagger}{\inshortcatvariables,\ccatvariableof{[\catorder]}{\midsymbol,0}},
        \reflectionofat{\onehotmapof{\zerotuple{[\catorder]}}}{\ccatvariableof{[\catorder]}{\midsymbol,0},\ccatvariableof{[\catorder]}{\midsymbol,1}},
        \unitaryat{\ccatvariableof{[\catorder]}{\midsymbol,1},\outshortcatvariables}
    }{\inshortcatvariables,\outshortcatvariables} \, .
\end{align*}


\begin{figure}[t]
    \begin{center}
        \begin{tikzpicture}[scale=0.35, thick] % , baseline = -3.5pt

%    \drawsmalltensor{0}{9}{\corelabelsize $\fbasis$}
%    \draw (1,9) -- (2,9);
%    \drawsmalltensor{3}{9}{\corelabelsize $\paulixsymbol$}
%    \draw (4,9) -- (5,9);
%    \drawsmalltensor{6}{9}{\corelabelsize $\hgate$}
%
%    %\draw (-1,5) rectangle (1,7);
%    \drawsmalltensor{0}{6}{\corelabelsize $\fbasis$}
%    \draw (1,6) -- (2,6);
%    \drawsmalltensor{0}{2}{\corelabelsize $\fbasis$}
%    \draw (1,2) -- (2,2);
%    \draw (2,1) rectangle (4,7);
%    \drawtextnode{3}{4}{\corelabelsize $\unitarysymbol$}
%    \draw (4,6) -- (12,6);
%    \draw (4,2) -- (12,2);
%
%    %\draw (2,1) rectangle (4,3);
%    %\drawtextnode{3}{2}{\corelabelsize $\hgate$}
%
%
%    \drawtextnode{-3}{9}{\corelabelsize $\headvariable$}
%    \draw[dashed] (-2,7.5) -- (30.5,7.5);
%    \drawtextnode{-3}{4}{\corelabelsize $\shortcatvariables$}
%    \drawtextnode{0}{4.5}{\corelabelsize $\vdots$}
%
%    \draw[dashed] (7.5,12) -- (7.5,-1) node[below] {\corelabelsize $\qstateofat{0}{\shortcatvariables} \otimes \hminusbasat{\headvariable}$};
%
%    \begin{scope}[shift={(4,0)}]
%        \draw (3,9) -- (4,9);
%        \draw (4,8) rectangle (7,10);
%        \node[anchor=center] (text) at (5.5,9) {\corelabelsize $\qcbencodingof{\exformula}$};
%        \draw (5,6) -- (5,8);
%        \drawvariabledot{5}{6}
%        \draw (6,2) -- (6,8);
%        \drawvariabledot{6}{2}
%    \end{scope}
%
%    \drawtextnode{9.5}{11}{\corelabelsize $\reflectionof{\badstatesymbol}$}
%
%    \draw[dashed] (11.5,11) -- (11.5,0.5);

    \drawtextnode{9.5}{4}{\corelabelsize $\inshortcatvariables$}
    \draw (10,6) -- (12,6);
    \drawtextnode{11.5}{4.25}{$\vdots$}
    \draw (10,2) -- (12,2);
    \draw (12,1) rectangle (14,7);
    \drawtextnode{13}{4}{\corelabelsize $\unitaryof{\dagger}$}
    \draw (14,6) -- (15,6);
    \drawsmalltensor{16}{6}{\corelabelsize $\paulixsymbol$}
    \draw (14,2) -- (15,2);
    \drawsmalltensor{16}{2}{\corelabelsize $\paulixsymbol$}
    \draw (17,2) -- (21,2);
    \draw (17,6) -- (18,6);
    \drawsmalltensor{19}{6}{\corelabelsize $\paulizsymbol$}
    \draw (18.5,5) -- (18.5,4);
    \drawvariabledot{18.5}{4}
    \node[rotate=-60] at (18.9,2.75) {$\cdots$};
    \draw (19.5,5) -- (19.5,2);
    \drawvariabledot{19.5}{2}

    \draw (20,6) -- (21,6);
    \drawsmalltensor{22}{6}{\corelabelsize $\paulixsymbol$}
    \draw (23,6) -- (24,6);
    \drawsmalltensor{22}{2}{\corelabelsize $\paulixsymbol$}
    \draw (23,2) -- (24,2);
    \draw (24,1) rectangle (26,7);
    \drawtextnode{25}{4}{\corelabelsize $\unitarysymbol$}
    \draw (26,6) -- (28,6);
    \draw (26,2) -- (28,2);
    \drawtextnode{27}{4.25}{$\vdots$}
    \drawtextnode{28.5}{4}{\corelabelsize $\outshortcatvariables$}
    %\draw (11,9)--(28,9);

%    \drawtextnode{19}{11}{\corelabelsize $\reflectionof{0}$}

    \draw[dashed] (23.5,0) -- (23.5,9.25) -- (14.5,9.25) -- (14.5,0) -- cycle;
    \drawtextnode{16}{8.5}{\corelabelsize $-\reflectionof{\onehotmapof{\zerotuple{[\catorder]}}}$}

    \draw[dashed] (20.5,1.25) -- (20.5,8.75) -- (17.5,8.75) -- (17.5,1.25) -- cycle;
    \drawtextnode{19}{8}{\corelabelsize $-\reflectionof{\onehotmapof{\onetuple{[\catorder]}}}$}

    \begin{scope}[shift={(-12,0)}]
        \drawtextnode{9.5}{4}{\corelabelsize $\inshortcatvariables$}
        \draw (10,6) -- (12,6);
        \drawtextnode{11.5}{4.25}{$\vdots$}
        \draw (10,2) -- (12,2);
        \draw (12,1) rectangle (14,7);
        \drawtextnode{13}{4}{\corelabelsize $-\reflectionof{\qstate}$}
        \draw (14,6) -- (15,6);

        \draw (14,6) -- (16,6);
        \drawtextnode{14.5}{4.25}{$\vdots$}
        \draw (14,2) -- (16,2);
        \drawtextnode{16.5}{4}{\corelabelsize $\outshortcatvariables$}

        \drawtextnode{19}{4}{${=}$}
    \end{scope}


    %\draw[dashed]  -- ;
    %\draw[dashed]  -- ; % node[below] {\corelabelsize $\qstateofat{1}{\shortcatvariables} \otimes \hminusbasat{\headvariable}$};


\end{tikzpicture}
    \end{center}
    \caption{
        Preparation of the negated reflection operator $-\reflectionof{\qstate}$ about a state $\qstate$, which is prepared by a unitary $\unitarysymbol$ acting on the ground state $\onehotmapofat{\zerotuple{[\catorder]}}{\shortcatvariables}$.
        The multiply controlled Pauli-Z is the negated reflection $-\reflectionof{\onetuple{[\catorder]}}$, and extended to the negated reflection $-\reflectionof{\onetuple{[\catorder]}}$ about the ground state by adding layers of single qubit Pauli-X.
        When concatenating with $\unitaryof{\dagger}$ and $\unitarysymbol$, the $-\reflectionof{\qstate}$ results.
    }
    \label{fig:reflectionZRealization}
\end{figure}

%% Construction based on oracle with ancilla head qubit
Alternatively we can use the phasekick technique (\lemref{lem:phaseKick}) when acting on an disentangled ancilla qubit prepared in the Hadamard basis state
\begin{align*}
    \hminusbasat{\headvariable}
    = \frac{1}{\sqrt{2}} \cdot \coloredmatrixof{1 \\ -1}{\headvariable} \, .
\end{align*}
We notice, that
\begin{align*}
    \contractionof{
        \qcbencodingofat{\exformula}{\cheadvariable{\outsymbol},\cheadvariable{\insymbol},\indexedshortcatvariables},
        \hminusbasat{\cheadvariable{\insymbol}}
    }{\cheadvariable{\outsymbol}}
    = \begin{cases}
          \hminusbasat{\cheadvariable{\insymbol}} & \ifspace \exformula(\shortcatindices) = 0\\
          -\hminusbasat{\cheadvariable{\insymbol}} & \ifspace \exformula(\shortcatindices) = 1
    \end{cases}
\end{align*}

Let us consider the span of one-hot encoded models of $\lnot\exformula$, namely
\begin{align*}
    \subspaceof{\lnot\exformula}
    \coloneqq \spanof{\onehotmapofat{\shortcatindices}{\shortcatvariables}\wcols\exformula(\shortcatindices)=0}
\end{align*}
and the span of one-hot encoded models
\begin{align*}
    \subspaceof{\exformula}
    \coloneqq\spanof{\onehotmapofat{\shortcatindices}{\shortcatvariables}\wcols\exformula(\shortcatindices)=1} \, .
\end{align*}
The reflection across $\subspaceof{\lnot\exformula}$ is then performed by (see \figref{fig:reflectionCCRealization})
\begin{align*}
    \reflectionofat{\subspaceof{\lnot\exformula}}{\outshortcatvariables,\inshortcatvariables}
    = \contractionof{
        \hminusbasat{\cheadvariable{\outsymbol}},
        \qcbencodingofat{\exformula}{\cheadvariable{\outsymbol},\cheadvariable{\insymbol},\shortcatvariables},
        \hminusbasat{\cheadvariable{\insymbol}},
        \identityat{\shortcatvariables,\outshortcatvariables,\inshortcatvariables}
    }{\outshortcatvariables,\inshortcatvariables}
\end{align*}
We further have
\begin{align*}
    \reflectionofat{\subspaceof{\exformula}}{\outshortcatvariables,\inshortcatvariables}
    = -\reflectionofat{\left(\subspaceof{\exformula}\right)^{\perp}}{\outshortcatvariables,\inshortcatvariables}
    = -\reflectionofat{\subspaceof{\lnot\exformula}}{\outshortcatvariables,\inshortcatvariables} \, .
\end{align*}
When preparing a disentangled ancilla qubit $\headvariable$ in the basis state $\hminusbasat{\headvariable}$, the oracle $\qcbencodingof{\exformula}$ is effectively a negated reflection across the span of one-hot encoded models of $\exformula$.

%This further coincides with the effective negative reflection across the span of one-hot encoded models
%\begin{align*}
%    \subspaceof{\exformula}
%    \coloneqq\spanof{\onehotmapofat{\shortcatindices}{\shortcatvariables}\wcols\exformula(\shortcatindices)=1} \, .
%\end{align*}

%We will investigate different effective schemes to realize $\reflectionof{\badstatesymbol}$ later.
%\begin{itemize}
%    \item When $\qstateof{\goodstatesymbol}=\onehotmapofat{\onetuple{[\seldim]}}{\avariables}\otimes $ and $\qstate$
%    \item Reflection across the hyperplane $\spanof{\onehotmapofat{\shortcatindices}{\shortcatvariables}\wcols\exformula(\shortcatindices)=0}$ effectively by $\qcbencodingof{\exformula}$ acting on an auxiliary disentangled qubit in $\hminusbasat{\headvariable}$
%\end{itemize}

\begin{figure}[t]
    \begin{center}
        \begin{tikzpicture}[scale=0.35, thick] % , baseline = -3.5pt


    %\drawtextnode{9.5}{4}{\corelabelsize $\inshortcatvariables$}

    \drawsmalltensor{10}{8}{\corelabelsize $\hminusbas$}
    \draw (11,8) -- (12,8);
    \draw (12,7) rectangle (15,9);
    \drawtextnode{13.5}{8}{\corelabelsize $\qcbencodingof{\formula}$}
    \draw (13,7) -- (13,6);
    \drawvariabledot{13}{6}
    \draw (14,7) -- (14,2);
    \drawvariabledot{14}{2}
    \draw (15,8) -- (16,8);
    \drawsmalltensor{17}{8}{\corelabelsize $\hminusbas$}

    \drawtextnode{10.5}{4}{\corelabelsize $\inshortcatvariables$}
    \draw (11,6) -- (16,6);
    \drawtextnode{12.5}{4.25}{$\vdots$}
    \draw (11,2) -- (16,2);
    \drawtextnode{16.5}{4}{\corelabelsize $\outshortcatvariables$}
    %\draw (12,1) rectangle (14,7);
    %\drawtextnode{13}{4}{\corelabelsize $\unitaryof{\dagger}$}
    %\draw (14,6) -- (15,6);

    \begin{scope}[shift={(-12,0)}]
        \drawtextnode{9.5}{4}{\corelabelsize $\inshortcatvariables$}
        \draw (10,6) -- (12,6);
        \drawtextnode{11.5}{4.25}{$\vdots$}
        \draw (10,2) -- (12,2);
        \draw (12,1) rectangle (14,7);
        \drawtextnode{13}{4}{\corelabelsize $-\reflectionof{\subspaceof{\formula}}$}
        \draw (14,6) -- (15,6);

        \draw (14,6) -- (16,6);
        \drawtextnode{14.5}{4.25}{$\vdots$}
        \draw (14,2) -- (16,2);
        \drawtextnode{16.5}{4}{\corelabelsize $\outshortcatvariables$}

        \drawtextnode{19}{4}{${=}$}
    \end{scope}


    %\draw[dashed]  -- ;
    %\draw[dashed]  -- ; % node[below] {\corelabelsize $\qstateofat{1}{\shortcatvariables} \otimes \hminusbasat{\headvariable}$};


\end{tikzpicture}
    \end{center}
    \caption{
        Realization of the reflection across the subspace of one-hot encoded models to a Boolean function $\formula$.
        An ancilla qubit is prepared in the Hadamard basis state $\hminusbas$ and serves as the head qubit of the oracle $\qcbencodingof{\formula}$.
        By the phase kickback effect, the qubit stays disentangled and the reflection $-\reflectionof{\subspaceof{\formula}}$ is effectively performed on the qubits $\shortcatvariables$.
    }
    \label{fig:reflectionCCRealization}
\end{figure}
\subsection{Moment Estimation}\label{sec:amplitudeEstimation}

As an example, we here investigate how to use the introduced quantum routines to measure the mean of a function given a probability distribution.
Given a distribution $\probwith$, the mean of a propositional formula $\formulaat{\shortcatvariables}$ is the contraction
\begin{align*}
    \meanparam
    = \contraction{\probat{\shortcatvariables},\formulaat{\shortcatvariables}} \, .
\end{align*}

\begin{figure}
    \begin{center}
        \begin{tikzpicture}[scale=0.35, thick] % , baseline = -3.5pt

    \drawsmalltensor{0}{9}{\corelabelsize $\fbasis$}
    \draw (1,9) -- (2,9);
    \drawsmalltensor{3}{9}{\corelabelsize $\paulixsymbol$}
    \draw (4,9) -- (5,9);
    \drawsmalltensor{6}{9}{\corelabelsize $\hgate$}

    %\draw (-1,5) rectangle (1,7);
    \drawsmalltensor{0}{6}{\corelabelsize $\fbasis$}
    \draw (1,6) -- (2,6);
    \drawsmalltensor{0}{2}{\corelabelsize $\fbasis$}
    \draw (1,2) -- (2,2);
    \draw (2,1) rectangle (4,7);
    \drawtextnode{3}{4}{\corelabelsize $\unitarysymbol$}
    \draw (4,6) -- (12,6);
    \draw (4,2) -- (12,2);

    %\draw (2,1) rectangle (4,3);
    %\drawtextnode{3}{2}{\corelabelsize $\hgate$}


    \drawtextnode{-3}{9}{\corelabelsize $\headvariable$}
    \draw[dashed] (-2,7.5) -- (30.5,7.5);
    \drawtextnode{-3}{4}{\corelabelsize $\shortcatvariables$}
    \drawtextnode{0}{4.5}{\corelabelsize $\vdots$}

    \draw[dashed] (7.5,12) -- (7.5,-1) node[below] {\corelabelsize $\qstateofat{0}{\shortcatvariables} \otimes \hminusbasat{\headvariable}$};

    \begin{scope}[shift={(4,0)}]
        \draw (3,9) -- (4,9);
        \draw (4,8) rectangle (7,10);
        \node[anchor=center] (text) at (5.5,9) {\corelabelsize $\qcbencodingof{\exformula}$};
        \draw (5,6) -- (5,8);
        \drawvariabledot{5}{6}
        \draw (6,2) -- (6,8);
        \drawvariabledot{6}{2}
    \end{scope}

    \drawtextnode{9.5}{11}{\corelabelsize $\reflectionof{\badstatesymbol}$}

    \draw[dashed] (11.5,11) -- (11.5,0.5);

    \draw (12,1) rectangle (14,7);
    \drawtextnode{13}{4}{\corelabelsize $\unitaryof{\dagger}$}
    \draw (14,6) -- (15,6);
    \drawsmalltensor{16}{6}{\corelabelsize $\paulixsymbol$}
    \draw (14,2) -- (15,2);
    \drawsmalltensor{16}{2}{\corelabelsize $\paulixsymbol$}
    \draw (17,2) -- (21,2);
    \draw (17,6) -- (18,6);
    \drawsmalltensor{19}{6}{\corelabelsize $\paulizsymbol$}
    \draw (18.5,5) -- (18.5,4);
    \drawvariabledot{18.5}{4}
    \node[rotate=-60] at (18.9,2.75) {$\cdots$};
    \draw (19.5,5) -- (19.5,2);
    \drawvariabledot{19.5}{2}

    \draw (20,6) -- (21,6);
    \drawsmalltensor{22}{6}{\corelabelsize $\paulixsymbol$}
    \draw (23,6) -- (24,6);
    \drawsmalltensor{22}{2}{\corelabelsize $\paulixsymbol$}
    \draw (23,2) -- (24,2);
    \draw (24,1) rectangle (26,7);
    \drawtextnode{25}{4}{\corelabelsize $\unitarysymbol$}
    \draw (26,6) -- (28,6);
    \draw (26,2) -- (28,2);

    \draw (11,9)--(28,9);

    \drawtextnode{19}{11}{\corelabelsize $\reflectionof{0}$}
    \draw[dashed] (26.5,12) -- (26.5,-1) node[below] {\corelabelsize $\qstateofat{1}{\shortcatvariables} \otimes \hminusbasat{\headvariable}$};


\end{tikzpicture}
    \end{center}
    \caption{Amplification scheme of the at the models of a Boolean function $\exformula$ supported coordinates of a quantum state.
    The ancilla variable $\headvariable$ is prepared as an disentagled Hadamard basis vector and stays disentangled during the following reflections.
    The oracle $\qcbencodingof{\formula}$ effectively performs the reflection $\reflectionof{\badstatesymbol}$ across the bad state supported at the models of $\formula$.}
    \label{fig:ancillaGrover}
\end{figure}


We assume that we have
\begin{itemize}
    \item a circuit $\unitarysymbol$ preparing a state with measurement distribution $\probwith$, i.e.
    \begin{align*}
        \qstateofat{0}{\outshortcatvariables}
        = \contractionof{\onehotmapofat{\zerotuple{[\catorder]}}{\inshortcatvariables},
            \unitaryat{\inshortcatvariables,\outshortcatvariables}
        }{\outshortcatvariables}
    \end{align*}
    and
    \begin{align*}
        \absof{\qstateofat{0}{\shortcatvariables}}^2
        = \probwith
    \end{align*}
    \item a \computationCircuit{} for $\formulawith$
\end{itemize}


\subsubsection{Moments as rotation angles}

We construct a Grover operator, which rotation angle determines the moment.
The moment can therefore be estimated based on the measurement of the operators eigenvalue phases.

We have
\begin{align*}
    %\meanparam
    \sum_{\shortcatindices\in[2]^{\catorder} \wcols \formula(\shortcatindices)=1}
    \absof{\qstateofat{0}{\indexedshortcatvariables}}^2
    = \contraction{\probat{\shortcatvariables},\formulaat{\shortcatvariables}}  = \meanparam \, .
\end{align*}
This insight motivates us to define the orthonormal quantum states
\begin{align*}
    \qstateofat{\goodstatesymbol}{\shortcatvariables}
    &= \frac{1}{\sqrt{\meanparam}} \sum_{\shortcatindices\in[2]^{\catorder} \wcols \formula(\shortcatindices)=1}
    \qstateofat{0}{\indexedshortcatvariables} \cdot \onehotmapofat{\shortcatindices}{\shortcatvariables} \\
    \qstateofat{\badstatesymbol}{\shortcatvariables}
    &= \frac{1}{\sqrt{1-\meanparam}} \sum_{\shortcatindices\in[2]^{\catorder} \wcols \formula(\shortcatindices)=0}
    \qstateofat{0}{\indexedshortcatvariables} \cdot \onehotmapofat{\shortcatindices}{\shortcatvariables}
\end{align*}
and split the prepared state into a weighted sum of them as
\begin{align*}
    \qstateofat{0}{\shortcatvariables}
    = \sqrt{\meanparam} \cdot \qstateofat{\goodstatesymbol}{\shortcatvariables}
    + \sqrt{1-\meanparam} \cdot \qstateofat{\badstatesymbol}{\shortcatvariables} \, .
\end{align*}

%% Comparison with quantum rejection sampling
Note that here the mean parameter $\meanparam$ is the analogon to the acceptance probability $\acceptanceprob$ in quantum rejection sampling.
To be more precise, when aiming to sample from the models of $\exformula$ using a rejection sampling scheme starting with the uniform distribution, the acceptance probability would be $\meanparam$.
Instead of amplifying this acceptance probability we here want to measure it to retrieve the mean parameter.

%% Preparation of the effective Grover operator
We prepare the reflection $-\reflectionof{0}$ in the Pauli-Z scheme and the reflection $\reflectionof{\badstatesymbol}$ in the oracle scheme (see \secref{sec:groverOperator}).
Their concatenation results in the negated Grover operator, which causes a phase offset by $\pi$.



%The Grover operator $\groverofat{\probtensor,\formula}{\shortcatvariables^{\insymbol},\shortcatvariables^{\outsymbol},\avariable^{\insymbol},\avariable^{\outsymbol}}$ amplifying $\formula$
%\begin{align*}
%    (\unitaryof{\textdagger}\circ\reflectionof{0}\circ \unitarysymbol) \otimes \qcbencodingof{\exformula}
%\end{align*}
%acting on the initial state
%\begin{align*}
%    \onehotmapofat{0}{\shortcatvariables} \otimes \sqrt{\frac{1}{2}}\cdot\coloredmatrixof{1\\-1}{\avariable} \, .
%\end{align*}

%\subsubsection{Effective reflection by the oracle}

%The oracle acts as
%\begin{align*}
%    \contractionof{
%    \qstateofat{0}{\shortcatvariables},\hminusbasat{\cheadvariable{\insymbol}},
%    \qcbencodingofat{\exformula}{\cheadvariable{\insymbol},\cheadvariable{\outsymbol}}}{\cheadvariable{\outsymbol},\shortcatvariables}
%    =\hminusbasat{\cheadvariable{\outsymbol}}
%    \otimes \left(-\sqrt{\meanparam}\qstateof{\goodstatesymbol}+\sqrt{1-\meanparam}\qstateof{\badstatesymbol}\right) \, ,
%\end{align*}
%and is therefore an effective reflection across $\qstateof{\badstatesymbol}$.
%We can exploit the usual sparsity mechanisms to prepare the oracle.

%The Grover operator $\groverofat{\qstateof{0},\formula}{\cheadvariable{\insymbol},\cheadvariable{\outsymbol},\inshortcatvariables,\outshortcatvariables}$ consists in (see \figref{fig:ancillaGrover})
%\begin{itemize}
%    \item Effective reflection across $\qstateof{\goodstatesymbol}$ by $\qcbencodingof{\formula}$.
%    \item Negated reflection across $\qstateof{0}$ by $(\unitaryof{\textdagger}\circ\reflectionof{\fbasis}\circ \unitarysymbol)$ where $\reflectionof{\fbasis}$ is the reflection across the ground state $\onehotmapofat{\zerotuple{[\catorder]}}{\shortcatvariables}$ leaving $\headvariable$ invariant.
%\end{itemize}

\subsubsection{Retrieving the moment from the phase}

We now estimate the eigenvalue phases of the Grover algorithm (see \secref{sec:groverAngleMeasurement}) to read of the moment.

Note that the projection of the state $\qstateof{0}$ onto the eigen space to the eigen value $1$ vanishes, and the corresponding trivial phase is not measured by the Quantum Phase Estimation procedure.

The phase is related to the rotation angle by (see \secref{sec:groverAngleMeasurement})
\begin{align*}
    \rotanglesymbol
    \approx \pm \frac{1}{\twopii} \sum_{\catenumerator\in[n]} \phi_{\catenumerator} 2^{\catorder-\catenumerator-1} \, .
\end{align*}
The rotation angle is further related to the mean $\meanparam$ as
\begin{align*}
    \rotanglesymbol = 2 \cdot \sin^{-1}\left(\sqrt{\meanparam}\right) \, .
\end{align*}
Through measuring the phase $\phasecore\in\{\pm\rotanglesymbol\}$ of one of these eigenvalues we can therefore retrieve the mean by
\begin{align*}
    \meanparam
    = \left(\sinof{\frac{\phasecore}{2}}\right)^2 \, .
\end{align*}

\subsubsection{Measuring the phase}

\todo{We apply the Quantum Phase Estimation (\secref{sec:quantumPhaseEstimation}) to measure the moment.}

We prepare auxiliary qubits $\avariables$ and perform for each $\selindexin$ $2^{\selindex}$ repetitions of $\groverof{\probtensor,\exformula}$ controlled on $\avariableof{\selindex}$.
The inverse Fourier transform on the auxiliary qubits is supported on the (2-adic representation of the) phases on the eigenvalues, i.e. $\pm\rotanglesymbol$, and can be detected by a computational basis measure of the ancilla variables.

% For the mean
When measuring the phase register as $\aindexof{[\seldim]}$ we compute a phase
\begin{align*}
    \phasecore
    =2\pi \cdot \sum_{\selindexin} 2^{-(\selindex+1)} \aindexof{\selindex} % 0.\phi_0\phi_1\ldots\phi_{\seldim-1}
\end{align*}
which is either $\alpha$ or $-\alpha$.
We retrieve the mean in both cases by
\begin{align*}
    \meanparam
    = \sin^2\left(\pi \cdot \sum_{\selindexin} 2^{-(\selindex+1)} \aindexof{\selindex}\right) \, .
\end{align*}

\todo{Computational resources: Square root improvement?}